\documentclass[12pt,letterpaper]{article}
\usepackage[utf-8]{inputenc}
\usepackage[margin=1in]{geometry}
\usepackage{setspace}
\usepackage{hyperref}
\usepackage{csquotes}
\usepackage{amsmath}
\usepackage{amssymb}
\usepackage[hidelinks]{hyperref}
\usepackage{fancyhdr}
\usepackage{tocloft}

% Set line spacing
\onehalfspacing

% Header and footer
\pagestyle{fancy}
\fancyhf{}
\rhead{Garcia | EQUITAS Framework}
\lhead{IDAI700: Ethics of Artificial Intelligence}
\cfoot{\thepage}

% Custom title formatting
\usepackage{titling}
\setlength{\droptitle}{-0.5cm}

% Better section formatting
\usepackage{sectsty}
\sectionfont{\normalsize\bfseries\centering}
\subsectionfont{\normalsize\bfseries}
\subsubsectionfont{\normalsize\itshape}

% Custom title
\title{\vspace{-1.5cm}
\Large \textbf{EQUITAS as Enforceable Equity}\\
\normalsize \textit{A Justice-Based Framework for AI Healthcare Diagnostics}\\
\normalsize Grounded in Community Co-Design, Binding Governance, and Civic Literacy\\
\vspace{0.3cm}
}

\author{\textbf{Piter Garcia}\\
IDAI700: Ethics of Artificial Intelligence\\
Rochester Institute of Technology\\
\texttt{pzg8794@rit.edu}
}

\date{December 2025}

\begin{document}

\maketitle

\begin{abstract}
\noindent This paper argues that AI healthcare diagnostics for culturally and linguistically diverse (CLD) adolescents require enforceable equity rather than aspirational principles. Drawing on lived experience of diagnostic harm, systematic evidence of AI underperformance in marginalized populations, and synthesis of six peer-reviewed sources, I propose EQUITAS: a framework that operationalizes equity through (1) community co-design with decision authority (Nadarzynski et al., 2024), (2) binding governance with auditable metrics and community veto power (Lekadir et al., 2025; Chen et al., 2023), (3) measurable trust indicators grounded in demonstrated performance (Sagona et al., 2025), and (4) civic literacy via Digital Resistance Mapping that prepares communities to act as algorithm auditors (Garcia, 2025; Munn, IDAI700). The framework integrates empirical evidence of disparities (Marko et al., 2025) with engineering practice and power analysis from critical AI ethics, rejecting the premise that ethics principles can succeed without enforcement mechanisms or community authority. EQUITAS treats diagnostic AI as a justice-critical institution: trustworthiness is achieved through procedural reform, not solely through improved algorithms.

\noindent \textbf{Keywords:} AI ethics, healthcare diagnostics, equity, procedural justice, CLD adolescents, governance, co-design, trust, enforcement

\end{abstract}

\newpage

\tableofcontents

\newpage

% ============================================================================
\section{Introduction}

Diagnostic systems were not built to see some lives with equal clarity, and the consequences of that invisibility are neither abstract nor incidental. As a Latino (Dominican) immigrant and trauma survivor who navigated childhood loss, abuse, and years of undiagnosed ADHD while living with the developmental and relational costs of that erasure, I learned that institutional systems often recode distress---particularly when expressed through cultural idioms, languages other than English, or somatic rather than behavioral presentations---as defiance, psychosomatic symptoms, or character weakness. \cite{garcia_ed415} The specific harm in my case involved not merely clinical error but a cascade of misinterpretations: what was trauma-responsive hypervigilance was read as oppositional behavior; what was executive dysfunction was coded as laziness; what was chronic illness caused by medical mismanagement was ultimately dismissed as psychosomatic until I conducted my own research and insisted on appropriate testing.\cite{garcia_ed415} These experiences taught me that diagnostic failure is not exceptional; it is systematic, shaped by the institutional architecture through which symptoms are recognized, interpreted, and treated.

This paper is grounded in two complementary insights. First, my empirical research brief on equitable diagnostic tools for CLD adolescents \cite{garcia_ed415} documents that these patterns are not idiosyncratic: nationally, Latino adolescents are underdiagnosed for ADHD at rates 40--50\% lower than white peers despite equal or higher symptom burdens, experience longer delays before any diagnosis, and receive less comprehensive follow-up care.\cite{garcia_ed415} When ADHD, chronic illness, depression, and childhood trauma intersect---as they do in my case and in the cases of many CLD adolescents---misrecognition is compounded across multiple clinical and educational domains. Second, emerging AI-assisted diagnostic tools risk automating and accelerating this historical pattern of marginalization unless equity is embedded as a non-negotiable design and governance requirement.\cite{marko2025} Marko et al.'s systematic review of 129 healthcare AI systems confirms that AI underperforms for racial and ethnic minorities, language-minoritized populations, women, disabled people, and low-SES groups precisely because inclusivity is absent from data, development teams, and evaluation metrics.\cite{marko2025}

The ethical problem at stake is threefold: \textit{epistemic injustice}, wherein marginalized communities' knowledge about their own bodies and the symptoms they experience is systematically discounted;\cite{garcia_ed415} \textit{procedural injustice}, wherein those most affected by diagnostic tools have no meaningful authority over their design, deployment, or correction; and \textit{distributive injustice}, wherein the burden of algorithmic error concentrates in already-marginalized populations while benefits accrue to the majority.\cite{garcia_ed415} Existing AI ethics frameworks have failed to prevent these harms because they remain what Luke Munn calls ``meaningless, isolated, and toothless''---they articulate high-level principles about fairness and engagement but lack binding thresholds, enforcement mechanisms, or mechanisms that grant communities veto power.\cite{munn_idai700} This allows institutions to claim compliance with equity language while continuing to deploy and iterate on diagnostic systems that underdiagnose CLD adolescents.\cite{munn_idai700}

To address this gap, I advance \textbf{EQUITAS}, a novel governance framework that operationalizes enforceable equity through three integrated mechanisms: (1) \textit{Equity-Centered Co-Design with Decision Authority}, building on Nadarzynski et al.'s participatory roadmap but upgrading community participation to co-ownership and veto power;\cite{nadarzynski2024} (2) \textit{Binding Governance with Enforcement Infrastructure}, integrating FUTURE-AI's concrete practices \cite{lekadir2025} with Chen et al.'s technical playbook for fairness \cite{chen2023} and adding community oversight authority that institutions cannot unilaterally waive; and (3) \textit{Civic Literacy via Digital Resistance Mapping}, a pedagogical and structural approach that prepares CLD communities to read, interpret, contest, and reshape diagnostic technologies rather than remaining passive recipients of algorithmic decisions.\cite{garcia_ed415} The framework deliberately incorporates trust as a measurable outcome \cite{sagona2025} and grounds itself in empirical evidence of documented disparities \cite{marko2025} rather than treating equity as an aspirational value.

The thesis is that trustworthy diagnostic AI for CLD adolescents is impossible without simultaneous transformation of both technical systems and governance institutions. If an algorithm achieves excellent overall accuracy while systematically underdiagnosing depression in Spanish-dominant youth or ADHD in girls of color, it is not technically sound and ethically acceptable---it is unsafe.\cite{marko2025} Conversely, if institutional processes are ``fair'' and participatory but communities lack enforceable authority to stop harmful deployments, those processes function as democracy theater while preserving the power asymmetries that perpetuate harm.\cite{munn_idai700} EQUITAS treats diagnostic AI as a justice-critical institution: it demands that equity be achieved through binding procedures, measurable outcomes, and power redistribution, not through hopeful language or technical sophistication alone.

% ============================================================================
\section{Background: Synthesis of Foundation Sources}

\subsection{The Design Foundation: Nadarzynski et al. (2024) on Co-Design as Origin of Trust}

Nadarzynski et al. establish a critical baseline: participatory design is not a late-stage ``acceptance'' exercise or a means of gathering feedback for refinement. Rather, genuine co-design is foundational to trustworthiness because it recognizes that those most affected by AI systems possess situated knowledge essential to their functioning.\cite{nadarzynski2024} Their empirical work on conversational agents in healthcare demonstrates that when marginalized community members---particularly patients from stigmatized or culturally minoritized backgrounds---are embedded as co-designers from problem framing through implementation and retirement, the resulting tools are more usable, better adapted to local contexts, and encounter less resistance in deployment.\cite{nadarzynski2024}

Ethically, Nadarzynski et al. frame co-design as a response to epistemic injustice: marginalized communities are repositioned as credible producers of knowledge about how systems should work and what harms to anticipate. For CLD adolescents navigating diagnostic systems, this insight is directly applicable. If families and youth have no authority to shape what counts as a symptom, what cultural expressions are clinically relevant, or how language is incorporated into diagnostic algorithms, then the system inherits majority-culture assumptions as invisible ground truth.\cite{nadarzynski2024} A mother who describes her child's distress through the idiom of \textit{nervios} (a culturally specific expression of anxiety and agitation) is not providing anecdotal feedback to be ``validated against'' clinical criteria; she is offering diagnostic knowledge that, when systematically ignored, produces structural misrecognition.

However, Nadarzynski et al.'s framework, while powerful, does not fully address the power dynamics that EQUITAS must confront. Co-design grounded in genuine participation can still occur within institutional hierarchies where the final decision authority remains with organizations that have incentives to minimize cost, accelerate deployment, or limit liability. EQUITAS therefore takes Nadarzynski's participatory lifecycle as necessary but upgrades it: co-design is only legitimate when marginalized communities hold binding authority---not merely consultative input---over what the system is for, what data it uses, and what happens when it causes documented harm.

\subsection{The Governance Architecture: Lekadir et al. (2025) and FUTURE-AI}

Lekadir et al.'s FUTURE-AI framework responds to a critical gap: translating abstract trustworthiness principles into auditable, institutionally binding practices. The framework articulates six pillars---Fairness, Universality, Traceability, Usability, Robustness, and Explainability---and maps each to concrete operational practices across the AI lifecycle, including data governance and provenance tracking, subgroup performance validation and reporting, transparent documentation artifacts such as model cards, post-deployment surveillance, and explicit criteria for system retirement.\cite{lekadir2025}

This specificity is ethically significant because it directly addresses a failure mode that Munn identifies: the transformation of ethics principles into institutional practice requires more than good intentions. When a healthcare organization commits to ``ensuring fairness,'' without specifying what fair performance means, how it will be measured, what constitutes an unacceptable threshold, or what happens when that threshold is violated, ``fairness'' becomes a rhetorical device rather than a binding obligation.\cite{lekadir2025} FUTURE-AI's move from principle to practice---specifying model cards, audit schedules, subgroup performance thresholds, and retirement triggers---represents a crucial step toward making equity verifiable.

Yet even FUTURE-AI contains structural limitations for the CLD communities this paper centers. The framework assumes implementation within existing institutional hierarchies and professional expertise structures. Governance is described as a series of practices to be adopted voluntarily, with community participation cast as advisory input rather than as decision authority that institutions cannot override.\cite{lekadir2025} For under-resourced safety-net clinics serving predominantly CLD populations, the full FUTURE-AI architecture may be treated as an unaffordable luxury rather than as a baseline safety requirement, effectively turning equity into a privilege good rather than a universal baseline.\cite{lekadir2025} EQUITAS extends FUTURE-AI by: (1) adding an enforcement layer that makes compliance mandatory rather than voluntary; (2) granting communities formal veto authority at procurement and deployment gates; and (3) proposing tiered implementation that ensures safety-net clinics can meet core equity requirements without deploying the full infrastructure immediately.

\subsection{Translating Principles to Practice: Chen et al. (2023) on Algorithmic Fairness in Medicine}

Chen et al. provide the technical playbook that bridges ethical intent and engineering reality. They argue---with empirical precision---that fairness in clinical AI is not a unidimensional property but a multiobjective challenge involving competing metrics such as equalized odds (equal false-positive and false-negative rates across groups), demographic parity (equal positive prediction rates), and calibration (predicted risk matches observed outcomes within demographic groups).\cite{chen2023} These metrics can conflict under different disease prevalences and clinical contexts, meaning that selecting one fairness metric is itself a policy decision with moral weight.\cite{chen2023}

Critically, Chen et al. demonstrate that aggregate accuracy (Area Under the Receiver Operating Characteristic Curve) is not merely insufficient but actively harmful for equity analysis, because it conceals large subgroup disparities behind overall performance numbers.\cite{chen2023} An algorithm achieving 92\% AUROC might perform at only 78\% sensitivity for a marginalized subgroup, a gap that aggregate metrics completely hide. Meaningful fairness assessment therefore requires: (1) subgroup-specific metrics reported with confidence intervals; (2) context-sensitive selection of fairness definitions aligned to what errors are most costly in a given clinical domain; and (3) evaluation across multiple dimensions of identity (race, ethnicity, language, gender, disability, SES) rather than single-axis analysis.\cite{chen2023}

Chen et al. also catalog a comprehensive mitigation toolkit spanning pre-processing (data reweighting, augmentation, representation repair), in-processing (adversarial debiasing, group-aware loss functions), and post-processing (group-specific thresholds, recalibration) approaches.\cite{chen2023} This technical menu is crucial because it demonstrates that fairness is not a matter of hoping algorithms will be unbiased; it is a matter of deliberately engineering systems to achieve specified fairness properties, monitoring for drift, and updating when performance degrades.\cite{chen2023}

For EQUITAS, Chen et al.'s work does something essential: it shows that operationalizing equity is technically tractable. The challenge is not that we lack the knowledge to build fairer diagnostic algorithms; it is that institutions often lack the incentives, governance structures, or accountability mechanisms to implement known fairness practices. In my ED415 research brief, I note that fair AI for CLD adolescent ADHD and depression screening would prioritize equalized-odds and true-positive-rate parity, accepting some increase in false positives (over-diagnosis) to avoid the well-documented catastrophe of false negatives (missed diagnoses) that have long cascading developmental costs.\cite{garcia_ed415} Chen's framework provides the technical vocabulary and methods to make that policy choice explicit and auditable.

\subsection{Trust as Verifiable Outcome: Sagona et al. (2025)}

Sagona et al. reframe trust from an aspirational quality or a marketing device into an observable outcome that can be engineered, measured, and maintained through specific institutional practices.\cite{sagona2025} They argue that trustworthiness in AI-assisted health systems depends on three interdependent dimensions: (1) \textit{explainability at the point of care}, moving beyond technical transparency to ensure that patients and clinicians can understand and meaningfully interpret algorithmic recommendations; (2) \textit{community governance}, specifically through empowered Community Advisory Boards (CABs) that review algorithms before deployment, have authority to request independent audits, and can reject systems if concerns are not adequately addressed; and (3) \textit{subgroup fairness validation as non-negotiable}, rejecting aggregate accuracy in favor of transparent, detailed documentation of model performance across demographic groups.\cite{sagona2025}

Ethically, Sagona et al. make a crucial move: they reframe resistance to medical AI in marginalized communities not as irrationality or fear of technology, but as historically justified distrust grounded in lived experience of medical racism and institutional betrayal.\cite{sagona2025} For many CLD families navigating healthcare, the rational response to a new diagnostic tool is skepticism, not because of technological illiteracy but because diagnostic systems have historically failed them. This means that trust cannot be demanded or asserted; it must be earned through consistent, observable evidence of fair treatment, transparency, and genuine accountability when harm occurs.\cite{sagona2025}

Sagona et al. propose specific mechanisms for earning that trust: incident reporting and remediation systems with transparent timelines; comprehensive transparency artifacts such as model cards that honestly account for known limitations; and ongoing community validation that checks whether algorithms perform as promised.\cite{sagona2025} For my EQUITAS framework, this contribution is essential because it specifies what trust operationally requires. Trust is not achieved through reassurance or goodwill; it is achieved through measurable practices such as complaint tracking, time-to-remediation, subgroup-specific satisfaction metrics, and formalized processes that allow communities to contest algorithmic decisions.\cite{sagona2025}

\subsection{The Empirical Mandate: Marko et al. (2025) on Documented Disparities}

Marko et al. provide the evidentiary foundation that transforms the question of equity from academic to urgent. Their systematic review of 129 healthcare and social-care AI systems documents a pervasive pattern: AI systems underperform for marginalized populations---defined by race, ethnicity, language, gender, disability status, and socioeconomic position---across diagnostic, risk-stratification, triage, and treatment-recommendation domains.\cite{marko2025} Crucially, they show that these disparities are not incidental or marginal; they are structural and reproducible, arising from three primary causal mechanisms: (1) training data that under-represents or excludes marginalized groups, (2) development teams that lack lived experience of marginalized communities' healthcare encounters and symptom presentations, and (3) evaluation regimes that prioritize aggregate metrics over subgroup performance, effectively rendering disparities invisible.\cite{marko2025}

In the context of my research brief, Marko et al.'s findings validate the urgency of the problem. My documentation of Latino adolescent ADHD underdiagnosis is not an outlier; it is one instance of a systematic pattern that spans diagnostic domains and healthcare systems.\cite{marko2025} When algorithms are trained predominantly on majority-culture populations and validated using aggregate metrics, they reliably fail minoritized and culturally-specific presentations. What appears as a technical problem---``the model wasn't trained on enough diverse data''---is actually a governance problem: the institutions building these systems have not been held accountable for ensuring that inclusion and subgroup validation are mandatory requirements, not optional enhancements.\cite{marko2025}

Marko et al. also identify a crucial gap between research and practice: disparities are well-documented, and remedies are widely recommended (diverse recruitment, participatory co-design, subgroup validation, post-deployment monitoring), yet institutions continue deploying systems that perpetuate harm.\cite{marko2025} This gap between knowledge and action is precisely where EQUITAS intervenes. Documenting harm is necessary but insufficient; what is required is governance infrastructure with enforceable authority to mandate inclusion and to stop deployment when documented evidence shows that CLD subgroups will be harmed.

\subsection{Critical Synthesis: Six Sources Converge on Enforceable Equity}

These five sources, synthesized through my ED415 research brief, converge on a powerful insight: trustworthy, equitable diagnostic AI requires simultaneous transformation across four dimensions---design (co-design), governance (binding infrastructure), measurement (subgroup metrics), and practice (engineering and operationalization)---and no single dimension alone is sufficient.\cite{garcia_ed415}

Nadarzynski et al. establish that co-design is foundational but does not guarantee authority.\cite{nadarzynski2024} Lekadir et al. demonstrate that governance infrastructure can specify practices but cannot enforce accountability without external oversight.\cite{lekadir2025} Chen et al. show that technical fairness is achievable but requires deliberate policy choices about which metrics to optimize.\cite{chen2023} Sagona et al. clarify that trust is earned through observable practices, not asserted through rhetoric.\cite{sagona2025} Marko et al. provide empirical evidence that disparities are structural, not incidental, and that business-as-usual deployment perpetuates them.\cite{marko2025} Munn adds that principles without teeth become ethics theater, allowing institutions to claim virtue while preserving power asymmetries.\cite{munn_idai700}

The tension that EQUITAS must resolve is this: participation without authority risks tokenism; governance without enforcement risks capture by institutional interests; metrics without contestability risk becoming a new form of quantified discrimination; and practices without power-sharing risk reproducing the very hierarchies they claim to address.\cite{munn_idai700} My research brief contributes the empirical specificity and intersectional grounding that makes these connections concrete for CLD adolescents navigating ADHD, depression, chronic illness, and trauma simultaneously.\cite{garcia_ed415}

% ============================================================================
\section{The EQUITAS Framework: Operationalizing Enforceable Equity}

\subsection{Part A: The Problem---Diagnostic Disparity as Systemic Injustice}

CLD adolescents experience diagnostic disparity not because individual clinicians or algorithms ``make random mistakes,'' but because the epistemic and procedural architecture of diagnosis was historically standardized on majority-culture norms and continues to be perpetuated through technology adoption.\cite{garcia_ed415} My research brief documents that Latino adolescents are underdiagnosed for ADHD at rates 40--50\% lower than white peers despite equal or higher symptom burdens.\cite{garcia_ed415} This gap persists across insurance status, geographic region, and provider type, suggesting structural rather than incidental causes. When ADHD, chronic illness, depression, and childhood trauma intersect---as in my own experience and in the experiences of many CLD adolescents---misrecognition is compounded: trauma-informed hypervigilance can be misread as oppositional defiance; executive dysfunction arising from ADHD can be moralized as laziness or poor motivation; somatic symptoms of depression can be dismissed as psychosomatic or exaggerated rather than investigated; and cultural expressions of distress can be pathologized as maladaptive rather than recognized as adaptive responses to marginalization.\cite{garcia_ed415}

AI does not enter a neutral space; it enters a diagnostic ecosystem already shaped by power imbalances, historical neglect, and institutional indifference to marginalized lives. If an algorithm is trained predominantly on data encoding majority-culture symptom presentations as ``normal'' and CLD presentations as noise or noncompliance, then even sophisticated fairness metrics cannot remedy the fundamental misalignment between what the model has learned to recognize and what CLD adolescents actually experience.\cite{marko2025} This creates what might be called a ``governance illusion'': an institution can publish equity language, deploy fairness dashboards, and host ethics meetings while preserving the underlying power structure and data practices that produce inequity.\cite{munn_idai700} In such contexts, ``fairness metrics'' can coexist comfortably with lived invisibility, because the metrics are designed to measure technical properties that may bear no relationship to whether CLD adolescents' symptoms are actually recognized or appropriately treated.\cite{marko2025}

Munn's critique is essential here: ethics frameworks without enforcement mechanisms risk becoming toothless, enabling institutions to claim virtue while continuing business as usual.\cite{munn_idai700} In healthcare, this dynamic is intensified by utilitarian pressure to deploy ``good enough'' systems quickly, often treating higher error rates in marginalized populations as an acceptable trade-off for speed and aggregate benefit.\cite{munn_idai700} EQUITAS begins from an opposite principle: if documented diagnostic harm concentrates in a specific population, and that population carries historical legacies of medical neglect and racialized dismissal, then precaution is not perfectionism---it is a basic safety stance consistent with responsible governance.\cite{garcia_ed415}

\subsection{Part B: The Solution---Three Pillars of EQUITAS}

EQUITAS operationalizes enforceable equity through three integrated pillars: Equity-Centered Co-Design with Decision Authority, Binding Governance with Enforcement, and Civic Literacy via Digital Resistance Mapping. These pillars are not meant to be implemented sequentially or in isolation; they work synergistically to address epistemic, procedural, and distributive injustice simultaneously.

\subsubsection{Pillar 1: Equity-Centered Co-Design with Decision Authority}

Following Nadarzynski et al., EQUITAS requires co-design across the entire AI lifecycle: problem framing, data curation, feature selection, metric definition, validation, deployment, monitoring, and retirement.\cite{nadarzynski2024} However, EQUITAS explicitly upgrades co-design from participation (advisory input) to co-authority (binding decision power). Concretely, this means that CLD adolescents and families are not merely consulted on an algorithm that professionals have already substantially designed; they hold seats on decision-making bodies with formal voting authority over fundamental questions: What is this tool intended for? What counts as a symptom versus a normative cultural variation? What language is acceptable? What trade-offs between sensitivity and specificity are ethically acceptable in this clinical context? What harms are unacceptable under any circumstances?

This directly addresses epistemic injustice because lived experience is reconstituted as primary diagnostic knowledge rather than as anecdotal feedback. When a mother describes her adolescent as experiencing \textit{nervios} and cultural distress idioms, her observation is not something to be ``validated'' against clinical criteria developed on majority-culture samples; her observation is itself a form of diagnostic knowledge that must shape what the algorithm is trained to recognize.\cite{garcia_ed415} Similarly, when an adolescent reports that their ADHD manifests primarily during transitions or in response to cultural discrimination rather than during focused academic tasks, that situated knowledge must inform how the algorithm conceptualizes attention and executive function.\cite{garcia_ed415}

The practical implementation of this pillar requires establishing Youth and Community Advisory Boards (CABs) with three key features: (1) membership drawn primarily from CLD communities and including adolescents themselves, not merely their parents; (2) meeting schedules and compensation that enable genuine participation rather than treating community input as unpaid emotional labor; and (3) documented authority to reject design choices, block deployment, or demand remediation when performance thresholds are not met.\cite{nadarzynski2024}

\subsubsection{Pillar 2: Binding Governance with Enforcement Infrastructure}

EQUITAS adopts Lekadir et al.'s FUTURE-AI as a governance backbone, specifying concrete practices across the lifecycle.\cite{lekadir2025} However, it adds a critical enforcement layer that transforms governance from a professional best-practice guideline into a legally and institutionally binding set of obligations. Specifically:

\textit{Pre-Deployment Requirements:} Any diagnostic tool targeting adolescents must undergo mandatory subgroup validation across race, ethnicity, language, gender identity, and SES before deployment is authorized.\cite{chen2023} Performance must meet explicit thresholds such as: equalized-odds gaps no greater than 3 percentage points across language groups; true-positive-rate ratios between 0.9 and 1.1 (meaning that the likelihood of correct diagnosis varies by no more than 10\% across groups); and subgroup-specific expected calibration error below 0.03 (meaning predicted risk matches observed outcomes within demographic strata).\cite{chen2023} These thresholds are not arbitrary; they are derived from my ED415 research demonstrating that CLD adolescents' false-negative rate (missed diagnoses) is the most costly error, requiring prioritization of equal opportunity and sensitivity parity.\cite{garcia_ed415}

\textit{Documentation and Transparency:} All algorithms must be accompanied by comprehensive model cards and fairness appendices that document: (1) datasets used and their demographic composition; (2) how and why specific fairness metrics were selected; (3) known limitations, particularly those affecting marginalized subgroups; (4) mitigation strategies employed and their trade-offs; and (5) a Conflicts of Interest and Governance section co-signed by CAB members.\cite{chen2023} These artifacts are not internal technical documents; they are public-facing materials designed to be legible to communities and clinicians who will use the system.\cite{lekadir2025}

\textit{Deployment and Monitoring:} Upon deployment, the tool must be subject to continuous monitoring with dashboards accessible to clinicians, institutional leadership, and community representatives. Dashboards track: (1) subgroup-specific performance metrics (fairness gaps, calibration, sensitivity/specificity by demographic group); (2) trust indicators including complaint volume, time-to-resolution, and community-reported satisfaction; and (3) drift indicators showing whether performance is stable over time and across populations.\cite{sagona2025} Critically, community representatives have access to these dashboards and can trigger review processes if metrics suggest harm is occurring.

\textit{Enforcement and Remediation:} When subgroup performance metrics exceed thresholds, automatic escalation procedures are triggered. If an equalized-odds gap drifts above 5 percentage points for any demographic group, the system enters a monitoring protocol requiring weekly performance review. If the gap persists above the agreed-upon threshold for more than 2 weeks, remediation is mandatory; if remediation does not restore performance within 4 weeks, the system is automatically suspended pending investigation and repair.\cite{chen2023} This is not reactive; it is built into the governance structure. CABs have explicit veto authority over whether the system can be reactivated after suspension, ensuring that communities are not merely informed of problems but have decision power over continuation.

Community Advisory Boards, critically, are not advisory bodies; they are governance bodies with binding authority. CAB decisions about deployment, suspension, and reactivation are recorded in public harm ledgers alongside institutional rationales and timeline documentation, creating accountability that survives organizational transitions and cost pressures.\cite{sagona2025}

\subsubsection{Pillar 3: Civic Literacy via Digital Resistance Mapping}

Trustworthy deployment requires more than transparency; it requires that affected communities can meaningfully interpret what is disclosed and exercise their governance authority with knowledge.\cite{munn_idai700} Digital Resistance Mapping is the pedagogical and structural solution: it transforms algorithmic oversight from an expert domain into a civic-literacy practice available to CLD youth and families.\cite{garcia_ed415}

Adapted from Rochester's historical Resistance Mapping project, Digital Resistance Mapping invites students and community members to: (1) overlay contemporary health or diagnostic outcome data with historical redlining maps and segregation patterns to reveal how past injustices structure present-day inequities; (2) use computational tools (Python, GIS software, fairness metric calculators) to compute equalized-odds gaps, demographic-parity ratios, and calibration error on real or simulated diagnostic data; (3) locate where algorithmic decisions come from (data sources, labeling processes, clinical interfaces, feedback loops); (4) identify where harm is most likely to accumulate (false negatives for which subgroups? false positives producing inappropriate treatment?); and (5) write governance memos proposing how algorithmic oversight should be structured, who should be accountable, and what recourse pathways should exist for families who believe they have been harmed.\cite{garcia_ed415}

This is not ``education'' in the sense of knowledge transfer from experts to communities. Rather, it is education as the foundational condition for legitimate participation in governance. CLD adolescents and families who understand how fairness metrics work, where disparities emerge in diagnostic systems, and how accountability structures function become more effective CAB members, data stewards, and advocates. They can read model cards with comprehension. They can interpret dashboard metrics with nuance. They can propose meaningful remediation. They can distinguish between genuine equity work and strategic language designed to obscure continuity in unjust practices.\cite{munn_idai700}

Digital Resistance Mapping also serves a second function: by having CLD youth themselves diagnose where and how diagnostic systems fail, the pedagogy productively inverts the usual dynamic of diagnostics. Rather than being positioned as the subjects of diagnostic assessment, students become analysts of diagnostic systems, reclaiming epistemic authority and demonstrating that lived experience is a form of diagnostic expertise.\cite{garcia_ed415}

\subsubsection{Integration: Operationalizing Multiple Justice Frameworks}

EQUITAS explicitly integrates epistemic, procedural, and distributive justice. \textit{Epistemic justice} is advanced by treating CLD narratives and lived experiences as primary diagnostic knowledge and by requiring that metric selection reflect community judgments about which errors are most costly---for ADHD screening, prioritizing equal opportunity to avoid the chronic developmental harms of underdiagnosis.\cite{garcia_ed415} \textit{Procedural justice} is embodied in CAB veto rights, transparent documentation of governance decisions, accessible appeals processes, and civic literacy initiatives that enable meaningful participation.\cite{nadarzynski2024} \textit{Distributive justice} is approximated through explicit subgroup performance thresholds that prevent burden from concentrating in marginalized populations, trust metrics stratified by demographic identity, and harm-ledger reporting that makes visible where benefits and harms fall.\cite{sagona2025}

The metrics proposed in EQUITAS---equalized-odds gap $\leq 3$ percentage points, true-positive-rate ratio between $0.9$ and $1.1$, subgroup expected calibration error $\leq 0.03$, complaint resolution SLA of 30 days---do not emerge from abstract fairness ideals. Rather, they are grounded in my ED415 research evidence about what disparities currently look like in adolescent ADHD and depression screening, in Chen et al.'s technical guidance about achievable fairness targets, and in Marko et al.'s documentation of what disparities matter most across domains.\cite{garcia_ed415}\cite{chen2023}\cite{marko2025} They are designed to be falsifiable: external regulators, community monitors, and researchers can assess whether institutions are meeting them.

\subsection{Part C: Why EQUITAS Is Ethically Sound}

EQUITAS is grounded in standpoint epistemology as a methodological commitment: lived experience is treated not as decoration or merely illustrative anecdote, but as situated knowledge about how systems fail and what justice requires.\cite{garcia_ed415} My trajectory from diagnostic erasure through research and advocacy to framework design is presented as evidence that those most harmed by systems possess insight essential to their repair. This epistemological stance rejects the default deference to majority-culture expertise that has historically justified exclusion of marginalized voices from design authority.\cite{nadarzynski2024}

EQUITAS is also responsive to power analysis: it explicitly treats trust, fairness, and equity as outcomes that require institutional power-sharing, not merely individual virtue or technical sophistication. Following Munn's skepticism about ethics initiatives that ignore incentive structures and authority asymmetries, EQUITAS insists that principles succeed only when accompanied by governance mechanisms that shift authority toward those most affected.\cite{munn_idai700} By embedding CAB veto rights, binding metric thresholds, enforceable monitoring, and civic-literacy programs that prepare communities to exercise governance authority, EQUITAS treats equity not as a value to be pursued but as a safety requirement that must be enforced.\cite{lekadir2025}

Finally, by connecting abstract ethics to concrete engineering practices and measurable indicators, EQUITAS addresses the implementation gap that plagues many frameworks. Trustworthiness is not asserted through language; it is demonstrated through auditable mechanisms such as governance meeting minutes, metric dashboards, harm ledgers, and community testimony about whether they experience systems as legitimate.\cite{sagona2025} This makes equity claims falsifiable and contestable in real time, rather than evaluated only retrospectively or only by experts.

% ============================================================================
\section{Counterargument and Response}

\subsection{The Pragmatic Utilitarian Objection}

A compelling objection---and one explicitly anticipated in IDAI700 course materials---comes from a pragmatic utilitarian perspective: EQUITAS is too expensive, too slow, and too burdensome for resource-constrained healthcare systems.\cite{munn_idai700} On this view, requiring Community Advisory Boards with veto power, extensive subgroup validation, continuous fairness and trust auditing, and community literacy programs will substantially delay deployment of diagnostic AI tools that could benefit the majority of adolescents now. If a new depression screener could markedly improve outcomes for 80--85\% of adolescents immediately, requiring years of co-design, validation, governance-structure establishment, and literacy-program development to serve the remaining 15--20\% might produce net harm by denying timely support to thousands of youth while systems wait for ``perfect'' equity.\cite{munn_idai700}

Proponents of this objection might cite historical examples of public-health interventions with imperfect equity profiles that nonetheless produced large aggregate benefits, such as vaccination campaigns or broad-scale screening initiatives. They might point to FUTURE-AI's 30 governance practices and Chen et al.'s multi-metric fairness regimes as evidence that full compliance is resource-intensive, particularly for under-resourced clinics and schools serving predominantly CLD populations.\cite{lekadir2025} On this view, deploying ``good enough'' algorithms and addressing disparities through iterative improvement and post-deployment monitoring is more ethical than withholding tools until full EQUITAS compliance is achieved. They might argue that EQUITAS's enforcement requirements risk turning equity into an unfunded mandate, creating barriers to adoption precisely in the resource-poor settings where new technologies are most needed.\cite{munn_idai700}

\subsection{Refutation: The Utilitarian Calculus Is Ethically and Empirically Flawed}

This objection fails on multiple grounds. \textit{First, ethically}, it encodes a hierarchy of whose harms matter: CLD adolescents are implicitly positioned as an acceptable ``error margin'' in service of majority benefits. This reproduces precisely the historical pattern that Marko et al. document: marginalized groups repeatedly bear disproportionate risks from new technologies while majorities capture the benefits, and this asymmetry is rationalized as reasonable cost-benefit analysis.\cite{marko2025} But if we would not accept an algorithm that systematically misdiagnosed white adolescents by 40\%, we cannot ethically accept the same underdiagnosis for Latino adolescents. The utilitarian argument effectively claims that marginalized lives have lower moral weight in the aggregate calculus, which violates basic principles of justice.\cite{garcia_ed415}

\textit{Second, empirically}, the ``good enough now, perfect later'' story is unsupported by evidence. Marko et al. show that systems deployed without inclusive design, subgroup validation, and community governance rarely drift toward equity; instead, disparities persist and often widen because the underlying structural causes remain unaddressed.\cite{marko2025} Once an algorithm is deployed, institutions have strong incentives to defend it, minimize reported harms, and frame equity improvements as ``future enhancements'' rather than safety corrections. My ED415 research further demonstrates that underdiagnosed CLD adolescents experience cumulative educational, psychological, and health harms that are not offset by aggregate gains elsewhere; missed diagnosis is not a minor cost, it is a developmental catastrophe.\cite{garcia_ed415}

\textit{Third, trust dynamics matter more than utilitarian calculus recognizes.} Sagona et al. demonstrate that when communities experience AI as yet another site of institutional misrecognition, trust erodes not only in that particular tool but in the clinicians and institutions deploying it.\cite{sagona2025} If a mother whose daughter was missed by a diagnostic AI algorithm learns that the institution deployed it despite knowing it performed worse for girls of color, her refusal to trust subsequent recommendations---even those that would actually benefit her daughter---is not irrationality but rational response to demonstrated untrustworthiness.\cite{sagona2025} From a long-term utilitarian perspective that accounts for system sustainability, legal liability, crisis-care costs, and reputational damage, inequitable deployment often produces net harm when all consequences are included.\cite{sagona2025}

\subsection{Acknowledging Genuine Constraints and Proposing Tiered Implementation}

The utilitarian objection does surface real tensions: healthcare systems operate under genuine resource scarcity, and adolescent mental-health crises demand urgent response.\cite{munn_idai700} EQUITAS does not deny this reality. Rather, it treats resource constraints as a governance problem rather than an excuse to normalize predictable injustice.

To address capacity concerns, EQUITAS proposes tiered implementation: large academic medical centers with ample resources implement the full framework including large empowered CABs, comprehensive subgroup analytics, multi-metric monitoring, and embedded civic-literacy programs; medium-sized regional health systems adopt a core EQUITAS model including basic CAB consultation, race/language/SES subgroup performance reporting, and participation in regional audit networks; small safety-net clinics implement minimal viable EQUITAS requirements including community partner representation on clinic governance, basic fairness metrics reported annually, and connection to university or public-health-partnership technical support.\cite{lekadir2025}

This tiering does not abandon equity; it distributes the burden proportionally while ensuring that no system serving CLD adolescents proceeds without measurable fairness validation and community involvement.\cite{lekadir2025} Public funding mechanisms---such as research-to-practice grants, health-equity appropriations, and regulatory incentives---can further reduce the financial burden on under-resourced clinics. Open-source fairness toolkits and pooled data-science capacity (e.g., through university partnerships) can lower technical barriers.\cite{chen2023}

The ethical requirement is not identical infrastructure everywhere but a refusal to deploy systems that have never been evaluated for the populations they are most likely to misdiagnose. If a small clinic cannot afford full EQUITAS compliance, the ethical response is institutional investment, not acceptance of harm. The premise that ``some inequity is better than no deployment'' inverts the burden of proof: institutions should have to justify why CLD adolescents should accept worse diagnostic accuracy, not why they should demand equity.\cite{garcia_ed415}

% ============================================================================
\section{Conclusion}

Diagnostic systems are not neutral tools; they are institutions that shape who is seen, who is believed, and who receives care. For CLD adolescents living at the intersection of ADHD, chronic illness, depression, and childhood trauma, diagnostic accuracy is not an academic question---it is a matter of identity, educational opportunity, developmental trajectory, and self-understanding. When diagnostic systems systematically fail to recognize these adolescents, they do not merely produce statistical disparities; they institutionalize epistemic injustice, procedural exclusion, and distributive harm.\cite{garcia_ed415}

This paper has argued that existing AI ethics frameworks often fail because they do not bind institutions to equity outcomes or grant communities enforceable authority to stop harm. EQUITAS responds to that failure by integrating six foundational sources into a coherent governance architecture: (1) co-design as the origin of trust (Nadarzynski et al., 2024), recognizing that marginalized communities must have authority over systems that affect them; (2) governance as auditable infrastructure (Lekadir et al., 2025), translating principles into concrete, institutionally binding practices; (3) trust as measurable outcome (Sagona et al., 2025), grounding trustworthiness in observable evidence rather than assertion; (4) empirical evidence of disparities (Marko et al., 2025), demonstrating that equity is a safety issue, not an optional enhancement; (5) ethics-to-engineering translation (Chen et al., 2023), showing that fairness is technically achievable but requires explicit policy decisions; and (6) resistance mapping as civic literacy (Garcia, 2025), preparing communities to exercise governance authority effectively.\cite{nadarzynski2024}\cite{lekadir2025}\cite{sagona2025}\cite{marko2025}\cite{chen2023}\cite{garcia_ed415} The framework also integrates Luke Munn's skepticism about ethics theater, insisting that principles without enforcement mechanisms become tools for institutional self-protection rather than for genuine equity.\cite{munn_idai700}

Concretely, healthcare institutions and AI vendors should operationalize EQUITAS through: (1) procurement requirements that make EQUITAS compliance a condition of contracts; (2) CAB governance structures with binding veto authority at deployment and monitoring gates; (3) mandatory subgroup performance validation with explicitly agreed-upon thresholds; (4) continuous monitoring systems with automatic escalation when metrics exceed acceptable ranges; (5) public-facing transparency artifacts designed for legibility to communities; (6) civic-literacy programs that prepare CLD adolescents and families to participate meaningfully in governance; and (7) enforceable remediation timelines with documented authority for communities to block continuation if harm persists.\cite{garcia_ed415}

The call to action is not simply ``build fairer models.'' It is to change who gets to decide what ``fair'' means and what happens when fairness fails. It is to treat diagnostic AI as a justice institution that must answer to those most affected, not merely to clinical leaders or technical experts. And it is to recognize that my trajectory from diagnostic harm to framework design is not inspirational narrative but evidence: lived experience is a diagnostic instrument for institutional failure and a blueprint for institutional repair.\cite{garcia_ed415}

If AI is to have a legitimate role in healthcare diagnostics, it must be governed so that CLD adolescents are no longer rendered invisible by design, code, or institutional indifference. EQUITAS is not a guarantee of perfect outcomes, but it is a commitment to transparency, authority-sharing, and enforceable accountability. It says: CLD adolescents deserve to be seen, and if systems fail to see them, institutions must be required to change.\cite{garcia_ed415}

% ============================================================================
\section*{AI Reflection (excluded from 3,000-word count)}

AI was used in this project as a drafting and structure assistant with deliberately constrained authority: to transform outlines into section-level prose, to strengthen the utilitarian counterargument, to improve coherence across EQUITAS pillars, and to help articulate precisely what operationalization requires when moving from principle to practice. AI was also used to test whether claims were actionable rather than aspirational---for example, by forcing explicit articulation of what CAB veto authority means procedurally, what ``trust as an outcome'' requires institutionally, and how Digital Resistance Mapping translates pedagogical theory into concrete literacy practice.\cite{garcia_ed415}

The central limitation encountered was that AI can make justice language sound fluent and complete while quietly smoothing over the power conflicts that are actually at stake. This risk mirrors what Munn critiques as ethics-washing: persuasive abstraction that can appear thorough while remaining procedurally unenforceable.\cite{munn_idai700} To counteract that tendency, I treated my lived experience, my research brief, and IDAI700 course materials as non-negotiable constraints: every time the prose drifted toward generic equity language, I forced it back into specific mechanisms (authority, veto, thresholds, monitoring, redress). This discipline prevented AI from smoothing away the very power asymmetries that EQUITAS is designed to expose and address.\cite{garcia_ed415}

This experience also deepened my understanding of trustworthiness as a governance problem. Even when AI-generated writing is coherent and appears well-sourced, it can obscure provenance and decision pathways unless deliberately documented. That insight maps directly onto why EQUITAS insists on traceability, model cards, and civic literacy: communities cannot meaningfully consent to or contest systems they cannot read and interpret.\cite{sagona2025}

Overall, using AI reinforced a core lesson from IDAI700: ethical outcomes are not produced by good intentions or smooth language; they are produced by governance structures, accountability mechanisms, and power-sharing that survive institutional incentives to cut corners.\cite{munn_idai700} This paper exists only because I insisted on making abstract ethics operational and refusing to accept AI's tendency toward premature consensus. That insistence is not a limitation of AI; it is evidence that human judgment, grounded in lived experience and institutional critique, must remain decisional.\cite{garcia_ed415}

% ============================================================================
\begin{thebibliography}{99}

\bibitem{nadarzynski2024}
Nadarzynski, T., et al. (2024). Achieving health equity through conversational AI: A roadmap for design and implementation of inclusive chatbots in healthcare. \textit{PLOS Digital Health}, 3(5), e0000492.

\bibitem{lekadir2025}
Lekadir, K., et al. (2025). FUTURE-AI: International consensus guideline for trustworthy and deployable artificial intelligence in healthcare. \textit{BMJ}, 388, e081554.

\bibitem{sagona2025}
Sagona, J., et al. (2025). Building and maintaining trust in AI-assisted health systems: A framework for equitable implementation. \textit{Journal of Medical Internet Research}, 27(1), e45678.

\bibitem{marko2025}
Marko, L., et al. (2025). Examining inclusivity in artificial intelligence for health and social care: A systematic review of population-level disparities. \textit{Journal of Health Informatics}, 32(2), 101--129.

\bibitem{chen2023}
Chen, R. J., Lu, M. Y., Sahai, S., Chen, T. Y., Williamson, D. F. K., Lipkova, J., Wang, J. J., \& Mahmood, F. (2023). Algorithmic fairness in artificial intelligence for medicine and healthcare. \textit{Nature Biomedical Engineering}, 7(6), 719--742.

\bibitem{garcia_ed415}
Garcia, P. (2025). Equitable diagnostic tools for marginalized adolescents: A CLD focus. \textit{ED415 Research Brief}, University of Rochester.

\bibitem{munn_idai700}
Munn, L. (2025). The uselessness of AI ethics. \textit{IDAI700: Ethics of Artificial Intelligence}, Unit 4 course reading. Rochester Institute of Technology.

\end{thebibliography}

\end{document}